The natural question is whether the $d$-logarithm factor is real or an artifact of the proof.
\begin{theorem}[Product Laplace forces the Gamma factor]
    \label{thm:product_laplace_gamma_lower}
    Fix $\lambda\in(0,1)$. For each $d\ge1$, let
    \[
        \ell(t)=\frac12 e^{-|t|},\qquad t\in\R,
        \qquad
        q_d(t)=\prod_{i=1}^d \ell(t_i),
        \qquad t=(t_1,\ldots,t_d)\in\R^d .
    \]
    Let $T_0,T_1,\ldots$ be i.i.d. samples from $q_d$, and define
    \[
        X=\lambda\sum_{k=0}^{\infty}(1-\lambda)^kT_k .
    \]
    Let
    \[
        P=P_{X,T_0},
        \qquad
        Q=P_X\otimes P_{T_0},
        \qquad
        R=\frac{dP}{dQ}.
    \]
    For $a>0$, define
    \[
        f_a(u)=e^{-a(u-1)}+a(u-1)-1,\qquad u\ge0,
    \]
    and
    \[
        D_{f_a}(P\|Q)=\E_Q[f_a(R)].
    \]
    Then there exist constants
    \[
        A_\lambda>0,\qquad
        \tau_\lambda<\infty,\qquad
        \beta_\lambda>0,
    \]
    depending only on $\lambda$, with $A_\lambda e^{\tau_\lambda}\ge2$, such that the following holds.
    With the universal constant $c_0=e^{-1}/8$, for every $d\ge1$ and every
    \[
        0<a\le (A_\lambda e^{\tau_\lambda})^{-d},
    \]
    we have
    \[
        D_{f_a}(P\|Q)
        \ge
        c_0
        \left(\frac{\beta_\lambda A_\lambda}{2}\right)^d
        a^2
        \sum_{j=0}^d
        \frac{L_a^j}{j!},
    \]
    where
    \[
        L_a
        =
        \log\left(
        \frac{1}{a(A_\lambda e^{\tau_\lambda})^d}
        \right).
    \]
    Consequently, for each fixed $\lambda\in(0,1)$ and $d\ge1$, there exist constants
    $c_{\lambda,d}>0$ and $a_{\lambda,d}\in(0,1)$ such that for all
    $0<a\le a_{\lambda,d}$,
    \[
        D_{f_a}(P\|Q)
        \ge
        c_{\lambda,d}
        a^2
        \sum_{j=0}^d
        \frac{\log^j(1/a)}{j!}.
    \]
    In particular,
    \[
        D_{f_a}(P\|Q)
        =
        \Omega_{\lambda,d}
        \left(
        a^2\frac{\log^d(1/a)}{d!}
        \right)
        \qquad\text{as }a\downarrow0.
    \]
\end{theorem}

\begin{proof}
    The proof isolates the source of the $d$ logarithms.  In one dimension,
    if the first Laplace sample is large and positive while the remaining
    contribution is bounded, then the likelihood ratio grows like $e^t$.
    In $d$ dimensions this effect tensorizes, so the likelihood ratio grows
    like $\exp(\sum_i t_i)$.  The Laplace density contributes
    $\exp(-\sum_i t_i)$, and truncating at height $1/a$ leaves exactly a
    Gamma-tail polynomial.

    \paragraph{Step 1: the one-dimensional bad event.}
    Let $L_0,L_1,\ldots$ be i.i.d. $\mathrm{Laplace}(0,1)$ random variables, and define
    \[
        X_\lambda
        =
        \lambda\sum_{k=0}^{\infty}(1-\lambda)^kL_k,
        \qquad
        X_\lambda'
        =
        \lambda\sum_{k=1}^{\infty}(1-\lambda)^{k-1}L_k .
    \]
    Let $p_\lambda$ be the density of $X_\lambda$. Then
    $X_\lambda'\stackrel{d}{=}X_\lambda$, $X_\lambda'\perp L_0$, and
    \[
        X_\lambda=\lambda L_0+(1-\lambda)X_\lambda'.
    \]
    Therefore, at a point of the form
    \[
        X_\lambda=\lambda t+(1-\lambda)z,
        \qquad
        L_0=t,
        \qquad
        X_\lambda'=z,
    \]
    the one-dimensional likelihood ratio is
    \[
        r_\lambda(t,z)
        =
        \frac{1}{1-\lambda}
        \frac{p_\lambda(z)}
        {p_\lambda(\lambda t+(1-\lambda)z)}.
        \label{eq:product_laplace_one_dim_ratio}
    \]

    We now prove that this ratio is exponentially large when $t$ is large
    and $z$ is bounded.  Write
    \[
        Z_\lambda
        =
        \lambda\sum_{k=1}^{\infty}(1-\lambda)^kL_k,
        \qquad
        X_\lambda=\lambda L_0+Z_\lambda.
    \]
    If
    \[
        q_\lambda(u)=\frac{1}{2\lambda}e^{-|u|/\lambda}
    \]
    is the density of $\lambda L_0$, then
    \[
        p_\lambda=q_\lambda*p_{Z_\lambda}.
    \]
    Define
    \[
        C_\lambda
        =
        \E\left[e^{|Z_\lambda|/\lambda}\right].
    \]
    This constant is finite, since
    \[
        \frac{|Z_\lambda|}{\lambda}
        \le
        \sum_{k=1}^{\infty}(1-\lambda)^k|L_k|,
    \]
    and hence, using $\E[e^{\alpha |L_0|}]=(1-\alpha)^{-1}$ for
    $0\le\alpha<1$,
    \[
        C_\lambda
        \le
        \prod_{k=1}^{\infty}
        \E\left[e^{(1-\lambda)^k|L_k|}\right]
        =
        \prod_{k=1}^{\infty}
        \frac{1}{1-(1-\lambda)^k}
        <\infty.
    \]
    Consequently, for every $x\ge0$,
    \[
        p_\lambda(x)
        =
        \int_\R
        \frac{1}{2\lambda}e^{-|x-z|/\lambda}
        p_{Z_\lambda}(z)\,dz
        \le
        \frac{C_\lambda}{2\lambda}e^{-x/\lambda}.
        \label{eq:product_laplace_density_tail}
    \]
    Since $p_\lambda$ is positive and continuous, also define
    \[
        m_\lambda
        =
        \inf_{|z|\le1}p_\lambda(z)
        >
        0.
    \]

    Now fix $|z|\le1$ and $t\ge(1-\lambda)/\lambda$. Then
    $\lambda t+(1-\lambda)z\ge0$, so \eqref{eq:product_laplace_density_tail} gives
    \[
        p_\lambda(\lambda t+(1-\lambda)z)
        \le
        \frac{C_\lambda}{2\lambda}
        \exp\left(
        -t-\frac{1-\lambda}{\lambda}z
        \right)
        \le
        \frac{C_\lambda}{2\lambda}
        \exp\left(
        -t+\frac{1-\lambda}{\lambda}
        \right).
    \]
    Substituting this into \eqref{eq:product_laplace_one_dim_ratio} yields
    \[
        r_\lambda(t,z)
        \ge
        \frac{2\lambda m_\lambda}{(1-\lambda)C_\lambda}
        \exp\left(-\frac{1-\lambda}{\lambda}\right)e^t .
    \]
    Set
    \[
        A_\lambda
        =
        \frac{2\lambda m_\lambda}{(1-\lambda)C_\lambda}
        \exp\left(-\frac{1-\lambda}{\lambda}\right).
    \]
    Choose $\tau_\lambda<\infty$ large enough that
    \[
        \tau_\lambda\ge\frac{1-\lambda}{\lambda}
        \qquad\text{and}\qquad
        A_\lambda e^{\tau_\lambda}\ge2.
    \]
    Then
    \[
        r_\lambda(t,z)\ge A_\lambda e^t,
        \qquad
        t\ge\tau_\lambda,\quad |z|\le1.
        \label{eq:product_laplace_one_dim_ratio_lower}
    \]
    Finally define
    \[
        \beta_\lambda
        =
        \PP(|X_\lambda'|\le1)
        =
        \int_{-1}^1 p_\lambda(z)\,dz
        >
        0.
    \]

    \paragraph{Step 2: tensorize the bad event.}
    Return to dimension $d$. Define
    \[
        X'
        =
        \lambda\sum_{k=1}^{\infty}(1-\lambda)^{k-1}T_k.
    \]
    Then $X'\stackrel{d}{=}X$, $X'\perp T_0$, and
    \[
        X=\lambda T_0+(1-\lambda)X'.
    \]
    Under $P=P_{X,T_0}$, write
    \[
        R_0=R(X,T_0)=\frac{dP}{dQ}(X,T_0).
    \]
    Since $q_d$ is a product distribution and the construction of $X$ is coordinatewise,
    the density of $X$ factorizes as
    \[
        p_X(x)=\prod_{i=1}^d p_\lambda(x_i).
    \]
    Therefore,
    \[
        R_0
        =
        \prod_{i=1}^d
        r_\lambda(T_{0,i},X_i').
        \label{eq:product_laplace_ratio_tensorization}
    \]

    Consider the event
    \[
        \mathcal G
        =
        \left\{
        |X_i'|\le1 \text{ and } T_{0,i}\ge\tau_\lambda
        \text{ for every } i=1,\ldots,d
        \right\}.
    \]
    On this event, \eqref{eq:product_laplace_one_dim_ratio_lower} and
    \eqref{eq:product_laplace_ratio_tensorization} imply
    \[
        R_0
        \ge
        A_\lambda^d
        \exp\left(\sum_{i=1}^d T_{0,i}\right).
        \label{eq:product_laplace_R0_large}
    \]
    Moreover, since $A_\lambda e^{\tau_\lambda}\ge2$, the same event implies
    \[
        R_0
        \ge
        (A_\lambda e^{\tau_\lambda})^d
        \ge2.
        \label{eq:product_laplace_R0_at_least_two}
    \]

    \paragraph{Step 3: turn the bad event into an $f_a$-divergence lower bound.}
    We first record an elementary lower bound on the generator.  For every
    $a\in(0,1]$ and every $u\ge2$,
    \[
        \frac{f_a(u)}{u}
        \ge
        c_0 a^2\min\left\{u,\frac1a\right\},
        \qquad
        c_0=\frac{e^{-1}}{8}.
        \label{eq:product_laplace_generator_lower}
    \]
    To prove this, let $s=a(u-1)$. If $s\le1$, then
    \[
        e^{-s}+s-1\ge \frac{e^{-1}}{2}s^2,
    \]
    and since $u\ge2$,
    \[
        \frac{(u-1)^2}{u}\ge \frac{u}{4}.
    \]
    Hence
    \[
        \frac{f_a(u)}{u}
        \ge
        \frac{e^{-1}}{8}a^2u
        \ge
        \frac{e^{-1}}{8}a^2\min\left\{u,\frac1a\right\}.
    \]
    If $s>1$, then it is enough to prove a lower bound of order $a$.
    When $1<s\le2$,
    \[
        f_a(u)=e^{-s}+s-1\ge e^{-1},
        \qquad
        u=\frac{s+a}{a}\le \frac{3}{a},
    \]
    so $f_a(u)/u\ge e^{-1}a/3$. When $s\ge2$,
    \[
        f_a(u)=e^{-s}+s-1\ge s-1\ge \frac{s}{2},
        \qquad
        u=\frac{s+a}{a}\le \frac{3s}{2a},
    \]
    so $f_a(u)/u\ge a/3$. This proves
    \eqref{eq:product_laplace_generator_lower}.

    Since $R=dP/dQ$, the change of measure identity gives
    \[
        D_{f_a}(P\|Q)
        =
        \E_Q[f_a(R)]
        =
        \E_P\left[\frac{f_a(R_0)}{R_0}\right].
    \]
    By the assumption $a\le(A_\lambda e^{\tau_\lambda})^{-d}$ and the fact
    $A_\lambda e^{\tau_\lambda}\ge2$, we have $a\le1$. Therefore, using
    \eqref{eq:product_laplace_generator_lower} on the event $\mathcal G$,
    where $R_0\ge2$, gives
    \[
        D_{f_a}(P\|Q)
        \ge
        c_0a^2
        \E_P\left[
            \min\left\{R_0,\frac1a\right\}
            \mathbf{1}_{\mathcal G}
            \right].
    \]
    Using \eqref{eq:product_laplace_R0_large}, and using the independence of
    $X'$ and $T_0$, we get
    \[
        \begin{aligned}
            D_{f_a}(P\|Q)
             & \ge
            c_0a^2\beta_\lambda^d
            \int_{[\tau_\lambda,\infty)^d}
            \min\left\{
            A_\lambda^d
            \exp\left(\sum_{i=1}^d t_i\right),
            \frac1a
            \right\}
            \prod_{i=1}^d
            \left(\frac12 e^{-t_i}\right)
            dt_i .
        \end{aligned}
        \label{eq:product_laplace_integral_lower}
    \]

    It remains only to evaluate the integral in
    \eqref{eq:product_laplace_integral_lower}. Set
    \[
        L_a
        =
        \log\left(
        \frac{1}{a(A_\lambda e^{\tau_\lambda})^d}
        \right).
    \]
    The assumption on $a$ gives $L_a\ge0$.  Put
    \[
        t_i=\tau_\lambda+u_i,
        \qquad
        u_i\ge0,
        \qquad
        S=\sum_{i=1}^d u_i.
    \]
    Then the integral in \eqref{eq:product_laplace_integral_lower} equals
    \[
        \begin{aligned}
            I_a
             & =
            2^{-d}
            \int_{\R_+^d}
            \min\left\{
            A_\lambda^d e^{d\tau_\lambda}e^S,
            \frac1a
            \right\}
            e^{-d\tau_\lambda}e^{-S}
            du
            \\
             & =
            2^{-d}
            \left[
                A_\lambda^d
                \operatorname{vol}\{u\in\R_+^d:S\le L_a\}
                +
                \frac{e^{-d\tau_\lambda}}{a}
                \int_{\{u\in\R_+^d:S>L_a\}} e^{-S}\,du
                \right].
        \end{aligned}
    \]
    The first term is a simplex volume:
    \[
        \operatorname{vol}\{u\in\R_+^d:S\le L_a\}
        =
        \frac{L_a^d}{d!}.
    \]
    For the second term, using the Gamma tail identity,
    \[
        \int_{\{u\in\R_+^d:S>L_a\}} e^{-S}\,du
        =
        \int_{L_a}^{\infty}
        e^{-s}\frac{s^{d-1}}{(d-1)!}\,ds
        =
        e^{-L_a}\sum_{j=0}^{d-1}\frac{L_a^j}{j!}.
    \]
    Since
    \[
        e^{-L_a}=a(A_\lambda e^{\tau_\lambda})^d,
    \]
    we obtain
    \[
        \frac{e^{-d\tau_\lambda}}{a}
        e^{-L_a}
        =
        A_\lambda^d.
    \]
    Hence
    \[
        I_a
        =
        2^{-d}A_\lambda^d
        \sum_{j=0}^d
        \frac{L_a^j}{j!}.
    \]
    Plugging this value of $I_a$ into
    \eqref{eq:product_laplace_integral_lower} gives
    \[
        D_{f_a}(P\|Q)
        \ge
        c_0
        \left(\frac{\beta_\lambda A_\lambda}{2}\right)^d
        a^2
        \sum_{j=0}^d
        \frac{L_a^j}{j!},
    \]
    which proves the main lower bound.

    \paragraph{Step 4: derive the asymptotic form.}
    For fixed $\lambda$ and $d$,
    \[
        L_a
        =
        \log(1/a)
        -
        d\log(A_\lambda e^{\tau_\lambda}).
    \]
    Therefore, if
    \[
        0<a\le
        a_{\lambda,d}
        \triangleq
        (A_\lambda e^{\tau_\lambda})^{-2d},
    \]
    then
    \[
        L_a\ge \frac12\log(1/a).
    \]
    Consequently,
    \[
        \sum_{j=0}^d\frac{L_a^j}{j!}
        \ge
        2^{-d}
        \sum_{j=0}^d
        \frac{\log^j(1/a)}{j!}.
    \]
    Absorbing the fixed factor
    \[
        c_0
        \left(\frac{\beta_\lambda A_\lambda}{2}\right)^d
        2^{-d}
    \]
    into $c_{\lambda,d}$ gives
    \[
        D_{f_a}(P\|Q)
        \ge
        c_{\lambda,d}
        a^2
        \sum_{j=0}^d
        \frac{\log^j(1/a)}{j!}
    \]
    for all $0<a\le a_{\lambda,d}$. The leading term in the sum is
    $\log^d(1/a)/d!$, which gives
    \[
        D_{f_a}(P\|Q)
        =
        \Omega_{\lambda,d}
        \left(
        a^2\frac{\log^d(1/a)}{d!}
        \right)
        \qquad\text{as }a\downarrow0.
    \]
\end{proof}